\section{Exact invariant-theory transports of the retained cubic and node}
\label{sec:gct-invariant-transport}

All schemes, coordinate rings and representations in this section are over
$\mathbb C$. The computations are defined over $\mathbb Q$, except for the
explicitly stated algebraic cube roots in the stabilizer. The original
inverse-fibre polynomial is retained as
\[
 P_{a,b,c}(r)=cr^3-2r^2+br-2a,
 \qquad P^h_{a,b,c}(R,T)=cR^3-2R^2T+bRT^2-2aT^3.
\]
In the source inverse chart, $r=y+1/x$, $\alpha=P'(r)/2$,
$x=\alpha^{-1}$, $y=r-\alpha$ and
$w=5\alpha^2-3r\alpha-c\alpha^3$; consequently the finite chart retains
only simple roots. The constructions below keep the algebraic completion
as a separate object. Their root at infinity when $c=0$ is a homogeneous
cubic root, not a licence to discard the source's additional point
$(0,b,a-4b^2)$.

The coordinate-ring restriction used in GCT is present explicitly in the
primary treatment of Ikenmeyer and Panova~\cite[Sections~1(a) and~2(b)]
{invariantIkenmeyerPanova}. Here we derive the particular
maps and representations directly. No determinant--permanent separation
is inferred from the following fixed-dimensional calculations.

\subsection{The binary representation, its coefficient slice and its discriminant}

Let $V$ be the four-dimensional space of binary cubics
$f=A R^3+B R^2T+C RT^2+D T^3$. Write $G_{\rm bin}=\operatorname{GL}_2(\mathbb C)$,
$z=(R,T)^{\mathsf T}$, and define the left action
\[
 (g\cdot f)(z)=f(g^{\mathsf T}z).
\]
Indeed $(g\cdot(h\cdot f))(z)=f(h^{\mathsf T}g^{\mathsf T}z)
=((gh)\cdot f)(z)$. For
$g=\left(\begin{smallmatrix}p&q\\s&t\end{smallmatrix}\right)$,
$\delta=pt-qs\ne0$, expansion gives
\begin{align*}
 A_g&=Ap^3+Bp^2q+Cpq^2+Dq^3,\\
 B_g&=3Ap^2s+B(p^2t+2psq)+C(2pqt+sq^2)+3Dq^2t,\\
 C_g&=3Aps^2+B(2pst+s^2q)+C(pt^2+2sqt)+3Dqt^2,\\
 D_g&=As^3+Bs^2t+Cst^2+Dt^3.
\end{align*}
The coefficient-slice embedding and its entire coordinate-ring map are
\[
 \iota:\mathbb A^3_{a,b,c}\longrightarrow V,
 \quad(a,b,c)\longmapsto(c,-2,b,-2a),
 \qquad
 \iota^*:\mathbb C[A,B,C,D]\longrightarrow\mathbb C[a,b,c],
\]
\[
 A\longmapsto c,\quad B\longmapsto-2,\quad C\longmapsto b,
 \quad D\longmapsto-2a.
\]
It is surjective: $a=-\iota^*D/2$, $b=\iota^*C$ and $c=\iota^*A$.
Its kernel is $(B+2)$, because evaluating a polynomial at $B=-2$
is division by the monic linear polynomial $B+2$, and the remaining
substitution $(A,C,D)=(c,b,-2a)$ is an invertible linear substitution.
The original affine slice itself is not $G_{\rm bin}$-stable.

Its saturation is the explicit equivariant morphism
\[
 \Phi:G_{\rm bin}\times\mathbb A^3\longrightarrow V,
 \qquad(g,a,b,c)\longmapsto g\cdot\iota(a,b,c),
\]
where $h$ acts on the source by $(g,a,b,c)\mapsto(hg,a,b,c)$.
Its pullback sends $A,B,C,D$ to the four displayed expressions with
$(A,B,C,D)=(c,-2,b,-2a)$, in
$\mathbb C[p,q,s,t,\delta^{-1},a,b,c]$. These formulas specify the
map without eliminating any root or selecting a branch.

\begin{proposition}\label{prop:invariant-saturation}
The image of $\Phi$ is $V\setminus\{0\}$, and $\ker\Phi^*=0$.
\end{proposition}
\begin{proof}
Every slice cubic has coefficient $-2$ and is nonzero; invertible
substitution cannot send it to zero. Conversely, take a nonzero cubic
$f$. There is a vector $z_0$ with $f(z_0)\ne0$: a polynomial vanishing
at every complex vector has all its coefficients zero, by successive
one-variable interpolation. Choose $w_0$ independent of $z_0$.
Euler's homogeneous identity, obtained by differentiating
$f(\lambda z_0)=\lambda^3f(z_0)$ at $\lambda=1$, gives
$df_{z_0}(z_0)=3f(z_0)\ne0$. Consequently
$df_{z_0}(w_0+kz_0)$ is nonzero for all but one $k\in\mathbb C$.
Choose one such $k$ and replace $w_0$ by $w=w_0+kz_0$.
The coefficient of $R^2T$ in $f(Rz_0+Tw)$ is exactly $df_{z_0}(w)$.
Replace $w$ by $-2w/df_{z_0}(w)$. The two chosen vectors remain
independent and the resulting coefficient is $-2$. Thus an explicitly
specified invertible variable substitution sends $f$ to the slice;
its inverse shows $f\in\operatorname{im}\Phi$. If a polynomial lies
in $\ker\Phi^*$, it vanishes on every nonzero vector of $V$ and hence
also at zero, by restricting to each line. Successive interpolation
again makes it the zero polynomial.
\end{proof}

Define, with no division by the leading coefficient,
\[
 \Delta(A,B,C,D)=B^2C^2-4AC^3-4B^3D-27A^2D^2+18ABCD.
\]
The substitution in $\iota$ gives exactly
\begin{equation}\label{eq:invariant-disc}
 \iota^*\Delta
 =4\bigl(b^2-cb^3+18abc-16a-27a^2c^2\bigr).
\end{equation}
For completeness, factor a binary cubic over $\mathbb C$ as
$f(z)=\prod_{i=1}^3(\alpha_iR+\beta_iT)$, allowing one factor to
carry any scalar. Direct expansion of the polynomial on the right
gives
\[
 \Delta(f)=\prod_{1\le i<j\le3}
              (\alpha_i\beta_j-\alpha_j\beta_i)^2.
\]
One may verify the identity first for
$f=A\prod_i(R-r_iT)$ by expanding the squared Vandermonde product:
it is $A^4\prod_{i<j}(r_i-r_j)^2$ and expands to the displayed
coefficient polynomial. Cubics with $A\ne0$ are dense, so equality
of these polynomials holds also at $A=0$. This proves that $\Delta=0$
is precisely a repeated projective root, including roots at infinity.
Under $g\cdot f$, the coefficient column of each linear factor is
multiplied by $g$. Each of the three two-by-two determinants is
multiplied by $\det g$, proving
\begin{equation}\label{eq:invariant-disc-weight}
 \Delta(g\cdot f)=(\det g)^6\Delta(f).
\end{equation}
The coordinate-ring representation is
$(g\cdot H)(f)=H(g^{-1}\cdot f)$, so the one-dimensional span of
$\Delta$ has character $(\det g)^{-6}$, not $(\det g)^6$.

\subsection{The actual source orbit and all of its stabilizer}

At the original target $(a,b,c)=(-1/4,0,0)$ put
\[
 q_0=-2R^2T+\tfrac12T^3,
 \qquad\Delta(q_0)=16.
\]
Its three projective roots are $[1:0]$, $[1/2:1]$ and $[-1/2:1]$.
We compute its stabilizer without identifying distinct scalar cubics.
Introduce the auxiliary polynomial $h(U,V)=UV(U-V)$ and the exact
invertible matrix
\[
 M=\begin{pmatrix}1&1/2\\0&1\end{pmatrix},\qquad
 M^{-1}=\begin{pmatrix}1&-1/2\\0&1\end{pmatrix},\qquad
 q_0(z)=-2h(Mz).
\]
This is an explicit comparison with an auxiliary polynomial; $q_0$
and its factor $-2$ remain fixed. Let
\begin{align*}
 D_1&=\begin{pmatrix}1&0\\0&1\end{pmatrix},&
 D_2&=\begin{pmatrix}-1&1\\0&1\end{pmatrix},&
 D_3&=\begin{pmatrix}0&1\\1&0\end{pmatrix},\\
 D_4&=\begin{pmatrix}1&0\\1&-1\end{pmatrix},&
 D_5&=\begin{pmatrix}0&1\\-1&1\end{pmatrix},&
 D_6&=\begin{pmatrix}1&-1\\1&0\end{pmatrix}.
\end{align*}
Substitution in the three factors of $h$ gives
\[
 h(D_jz)=\varepsilon_jh(z),\qquad
 (\varepsilon_1,\ldots,\varepsilon_6)=(1,1,-1,1,-1,-1).
\]
Their fractional-linear maps are respectively
$t$, $1-t$, $1/t$, $t/(t-1)$, $1/(1-t)$ and $(t-1)/t$.
They give all permutations of $\{0,1,\infty\}$: inspection gives
six distinct permutations, and a projective transformation fixing all
three must be the identity. To prove the latter assertion, fixing
$0$ and $\infty$ makes its matrix diagonal up to scalar; fixing $1$
makes the two diagonal entries equal.

It follows that the full stabilizer $H_0$ is the following set of
eighteen explicitly determined matrices:
\begin{equation}\label{eq:invariant-stab}
 H_0=
 \left\{\left(M^{-1}\kappa D_jM\right)^{\mathsf T}:
      1\le j\le6,\quad\kappa^3\varepsilon_j=1\right\}.
\end{equation}
Indeed a stabilizer must permute the three roots. After conjugating
its transpose by $M$, it is therefore $\kappa D_j$ for exactly one
$j$ and a nonzero scalar $\kappa$. Substitution into $-2h(Mz)$ makes
the multiplier exactly $\kappa^3\varepsilon_j$, which must be one.
Conversely every listed matrix has that multiplier and stabilizes
$q_0$. Each of the six scalar equations has three different roots,
and the six projective classes are different. The homomorphism to
root permutations is explicitly $[z]\mapsto[g^{-\mathsf T}z]$:
$(gh)^{-\mathsf T}=g^{-\mathsf T}h^{-\mathsf T}$ proves its composition
law. The transposes used to compute stabilizer membership above
give its inverse permutations. Consequently the root action gives
the exact sequence
\[
 1\longrightarrow\{\omega I:\omega^3=1\}
 \longrightarrow H_0\longrightarrow\mathfrak S_3\longrightarrow1.
\]

Every two unordered triples of distinct projective points are related
by a projective transformation. Here is a constructive proof. Choose
nonzero representatives $v_1,v_2$ for the first two points; they form
a basis. Write a representative for the third as $av_1+bv_2$;
distinctness gives $a,b\ne0$. The linear isomorphism sending
$av_1$ and $bv_2$ to $(1,0)^{\mathsf T}$ and $(0,1)^{\mathsf T}$
sends the third point to $[1:1]$. Apply this construction twice and
compose one map with the other's inverse. Consequently a squarefree
binary cubic is an invertible substitution of $q_0$ times a nonzero
scalar. Absorb that scalar by a scalar matrix with its specified
cube root. We have proved
\begin{equation}\label{eq:invariant-open-orbit}
 G_{\rm bin}q_0=\{f:\Delta(f)\ne0\},\qquad
 \overline{G_{\rm bin}q_0}=V.
\end{equation}
The closure equality also follows directly because the complement
is the zero set of a nonzero polynomial. Its coordinate ring and
the coordinate ring of the orbit are, respectively,
\[
 \mathbb C[\overline{G_{\rm bin}q_0}]=\mathbb C[A,B,C,D],
 \qquad
 \mathbb C[G_{\rm bin}q_0]=\mathbb C[A,B,C,D,\Delta^{-1}].
\]
The localization description is the defining coordinate ring of the
principal affine open set: it is represented by the hypersurface
$z\Delta-1=0$, with projection inverse $f\mapsto(f,\Delta(f)^{-1})$.

We can compute the invariant functions as well. Scalar matrices act
on $V$ by $f\mapsto\lambda^3f$. If a polynomial is $G_{\rm bin}$-invariant,
decompose it into its homogeneous terms. Comparing the distinct powers
$\lambda^{3d}$ in its value on $\lambda^3f$ makes every positive-degree
term zero. Hence $\mathbb C[V]^{G_{\rm bin}}=\mathbb C$.

In contrast,
\begin{equation}\label{eq:invariant-SL2}
 \mathbb C[V]^{\operatorname{SL}_2}=\mathbb C[\Delta].
\end{equation}
To prove this equality explicitly, first note that
$\det(H_0)=\{\zeta:\zeta^6=1\}$. Formula
\eqref{eq:invariant-disc-weight} gives containment. Scalar stabilizers
give all cube roots of unity as determinants, and the matrix with
$D_2$ and $\kappa=1$ gives determinant $-1$, so the reverse containment
holds. If $f_1$ and $f_2$ are squarefree and have equal discriminant,
choose $g$ with $g f_1=f_2$. Then $(\det g)^6=1$. The stabilizer of
$f_1$ is conjugate to $H_0$, so it has an element $h$ with
$\det h=(\det g)^{-1}$. Now $gh\in\operatorname{SL}_2$ still sends
$f_1$ to $f_2$. For an $\operatorname{SL}_2$-invariant polynomial $J$,
define the polynomial in one variable
\[
 j(t)=J\left(-2R^2T+\frac{t}{32}T^3\right).
\]
The argument of $J$ has discriminant $t$. Thus the just-proved orbit
statement gives $J(f)=j(\Delta(f))$ on $\Delta\ne0$, hence everywhere
by polynomial density. Conversely \eqref{eq:invariant-disc-weight}
makes every polynomial in $\Delta$ an $\operatorname{SL}_2$-invariant.

For later representation comparisons one need not suppress any weights.
For $g=\operatorname{diag}(\lambda,\mu)$ the full degree-$d$ character
of the coordinate ring of this orbit closure is
\begin{equation}\label{eq:invariant-cubic-character}
 \sum_{i+j+k+\ell=d}
 \lambda^{-(3i+2j+k)}\mu^{-(j+2k+3\ell)}.
\end{equation}
This follows from its monomial basis $A^iB^jC^kD^\ell$ and the four
coefficient weights. In particular every degree, every weight and
every multiplicity in this orbit closure agrees with the corresponding
ambient coordinate-ring representation, because the two rings are
equal by \eqref{eq:invariant-open-orbit}.

\subsection{The retained node as a ring and a representation}

On $b=c=0$, the source root cover is $a=-r^2$. Its completed
self-fibre product and its exact linear comparison are
\[
 B_{\rm root}=\mathbb C[r_1,r_2]/(r_1^2-r_2^2),\qquad
 R_{\rm node}=\mathbb C[u,v]/(uv),
\]
\[
 u=r_1-r_2,\quad v=r_1+r_2,\qquad
 r_1=(u+v)/2,\quad r_2=(v-u)/2.
\]
Substitution gives $uv=r_1^2-r_2^2$, and the displayed inverse proves
the ring isomorphism. In this ring,
\begin{equation}\label{eq:invariant-node-base}
 a=-r_1^2=-r_2^2=-\frac{u^2+v^2}{4}.
\end{equation}
Removing $r_1r_2=0$ removes the origin from both components, so the
crossing and the modules supported there belong to the completed object.

Let $K\subset G_{\rm node}=\operatorname{GL}_2(\mathbb C)$ be the monomial
subgroup generated by the diagonal torus
$T_{\rm node}=\{\operatorname{diag}(\lambda,\mu):\lambda\mu\ne0\}$
and $\sigma=\left(\begin{smallmatrix}0&1\\1&0\end{smallmatrix}\right)$.
It acts on node points by ordinary matrix multiplication, and on
functions by $(k\cdot f)(z)=f(k^{-1}z)$. Thus $u$ and $v$ have
torus characters $\lambda^{-1}$ and $\mu^{-1}$ and $\sigma$ exchanges
them. The relations are preserved. This $K$ does not fix the map to
$a$: a diagonal point transformation sends the expression in
\eqref{eq:invariant-node-base} to
$-(\lambda^2u^2+\mu^2v^2)/4$. All uses of $K$ below concern this
specified node action, not an unstated automorphism over $\mathbb A^1_a$.

For clarity retain $R_{\rm node}$ and define its comparison ring
\[
 \widetilde R=\mathbb C[u]\oplus\mathbb C[v],\qquad
 j(f)=(f(u,0),f(0,v)).
\]
Every element of $R_{\rm node}$ has a unique expression
$c+u p(u)+v q(v)$: the relation kills all mixed monomials, and
restriction to the two axes proves uniqueness. Thus $j$ is injective,
its image is exactly the pairs with equal constant terms, and the
inverse on that image is $(p(u),q(v))\mapsto p(u)+q(v)-p(0)$.
Every element of $\widetilde R$ is integral over the image, since
$\widetilde R$ is generated as an $R_{\rm node}$-module by $(1,1)$
and the idempotent $(1,0)$, which satisfies $e^2-e=0$.
The total ring of fractions is $\mathbb C(u)\oplus\mathbb C(v)$.
To check this assertion, the non-zero-divisor $u+v$ allows the
fractions $u/(u+v)=(1,0)$ and $v/(u+v)=(0,1)$; localizing within
each component gives its rational-function field. Each polynomial
ring in one variable is integrally closed: a rational function
$p/q$ with coprime $p,q$ satisfying a monic equation makes $q$
divide $p^{n}$, so $q$ is constant. Componentwise application proves
$\widetilde R$ integrally closed. It is therefore precisely the
algebraic normalization comparison ring, with the injection retained.

Define $\epsilon:\widetilde R\to\mathbb C$ by
$\epsilon(p,q)=p(0)-q(0)$. Its codomain is the $K$-module
$\mathbb C_{\rm sgn}$ on which the torus acts trivially and
$\sigma$ acts by $-1$. The preceding reconstruction proves the exact
$K$-equivariant sequence
\begin{equation}\label{eq:invariant-node-sequence}
 0\longrightarrow R_{\rm node}\xrightarrow{j}\widetilde R
 \xrightarrow{\epsilon}\mathbb C_{\rm sgn}\longrightarrow0.
\end{equation}
The conductor is precisely
\[
 \mathfrak c=(u,v),\qquad
 j(\mathfrak c)=u\mathbb C[u]\oplus v\mathbb C[v].
\]
Indeed a zero-constant pair remains an equal-constant pair after
multiplication by any element of $\widetilde R$. Conversely multiplying
a conductor element by $(1,0)$ forces its common constant to be zero.
The quotient in \eqref{eq:invariant-node-sequence} is annihilated by
$\mathfrak c$, and conversely a nonzero constant cannot annihilate it.

Give $u,v$ degree one. The degree-zero pieces of
\eqref{eq:invariant-node-sequence} are respectively the trivial
representation, the sum of trivial and sign representations, and the
sign representation. For every $n\ge1$ the two rings have the same
degree-$n$ basis $u^n,v^n$, with torus character
$\lambda^{-n}+\mu^{-n}$ and with $\sigma$ swapping its basis.
Therefore
\[
 R_{\rm node}^{T_{\rm node}}=\mathbb C,\quad
 \widetilde R^{T_{\rm node}}=\mathbb C\oplus\mathbb C,\quad
 R_{\rm node}^{K}=\widetilde R^K=\mathbb C.
\]
These equalities follow by comparing the independent torus characters
in the monomial bases and then taking the swap invariants. Thus the
full $K$-invariant scalar rings coincide, while the exact sign
cokernel in \eqref{eq:invariant-node-sequence} remains nonzero.

The differential module and its comparison map are
\[
 \Omega^1_{R_{\rm node}}=
 (R_{\rm node}\,du\oplus R_{\rm node}\,dv)/
       R_{\rm node}(v\,du+u\,dv),
 \qquad
 \Omega^1_{\widetilde R}=\mathbb C[u]du\oplus\mathbb C[v]dv.
\]
The presentation follows directly from the universal derivation of
the polynomial ring and $d(uv)=vdu+udv$. Its class
$\theta=vdu=-udv$ is nonzero. Indeed if $(v,0)=h(v,u)$ in the free
module, then $hu=0$, so the unique-form description gives $h\in(v)$.
The equation $hv=v$ would then imply $h(0,v)=1$, contradicting its
zero constant. Multiplying the relation by $u$ and $v$ proves
$u\theta=v\theta=0$. The non-zero-divisor $u+v$ kills $\theta$,
so this is torsion in the literal module-theoretic sense.

If $Adu+Bdv$ maps to zero in $\Omega^1_{\widetilde R}$, then
$A(u,0)=0$ and $B(0,v)=0$, giving $A=va(v)$ and $B=ub(u)$.
The relations $v^2du=0$ and $u^2dv=0$ reduce its class exactly to
$(a(0)-b(0))\theta$. Conversely $\theta$ maps to zero, and every
element of $\Omega^1_{\widetilde R}$ lifts by these two displayed
summands. We have proved the exact sequence
\begin{equation}\label{eq:invariant-differential-sequence}
 0\longrightarrow\mathbb C\theta\longrightarrow
 \Omega^1_{R_{\rm node}}\longrightarrow\Omega^1_{\widetilde R}
 \longrightarrow0.
\end{equation}
The original-coordinate identity is
$\theta=r_2dr_1-r_1dr_2$: substitute the two linear inverse formulas
and use $r_1dr_1=r_2dr_2$, the derivative of the original relation.

Differentiation commutes with the specified action on functions.
It gives
\[
 \operatorname{diag}(\lambda,\mu)\cdot\theta
 =\lambda^{-1}\mu^{-1}\theta,
 \qquad \sigma\cdot\theta=u\,dv=-\theta.
\]
Thus $\mathbb C\theta=\mathbb C_{\det^{-1}}$ as a $K$-module.
With $\deg du=\deg dv=1$, it lies in degree two, and
\[
 \operatorname{Hilb}(\Omega^1_{R_{\rm node}},t)
 =\frac{2t}{1-t}+t^2.
\]
More precisely, the degree-$n$ character for $n\ge1$ is
$\lambda^{-n}+\mu^{-n}$ plus the additional character
$\lambda^{-1}\mu^{-1}$ exactly at $n=2$; the additional swap
character is negative. No copy of $\det^{-1}$ occurs in
$R_{\rm node}$ itself, since none of its torus weights is
$\lambda^{-1}\mu^{-1}$. The exact module comparison
\eqref{eq:invariant-differential-sequence}, rather than replacement
of the module by its coefficient ring, retains this difference.

\subsection{An exact common torus and the discriminant pullback}

Define the two homomorphisms from $\mathbb G_m$ by
\[
 \beta(\eta)=\operatorname{diag}(\eta^{-1},\eta^2)\in G_{\rm bin},
 \qquad
 \nu(\eta)=\eta^3 I\in G_{\rm node}.
\]
The map from the completed node to binary cubics is
\begin{equation}\label{eq:invariant-node-to-cubic}
 \psi(u,v)=-2R^2T+\frac{u^2+v^2}{2}T^3
          =P^h_{a,0,0}(R,T),\quad a=-\frac{u^2+v^2}{4}.
\end{equation}
Its coordinate-ring map is
\[
 \psi^*: A\mapsto0,\ B\mapsto-2,\ C\mapsto0,
 \ D\mapsto(u^2+v^2)/2.
\]
Its kernel is exactly $(A,C,B+2)$. To see the unasserted part of
that equality, after division by those three generators any kernel
element would be a polynomial $H(D)$ with $H(u^2/2)=0$ on $v=0$;
its distinct powers of $u$ force every coefficient to vanish.
Its image is the polynomial subalgebra
$\mathbb C[(u^2+v^2)/2]$, not the entire node ring.
Since $\beta(\eta)$ multiplies $R^2T$ by
$\eta^{-2}\eta^2=1$ and $T^3$ by $\eta^6$, while $\nu(\eta)$
multiplies $u^2+v^2$ by $\eta^6$, this map is equivariant for the
two stated homomorphisms. Its discriminant pullback is exactly
\begin{equation}\label{eq:invariant-disc-node}
 \psi^*\Delta=16(u^2+v^2)=-64a.
\end{equation}
Both $\Delta$ restricted along $\beta$ and $\theta$ restricted along
$\nu$ have coordinate-module weight $\eta^{-6}$; their domains and
module structures remain different and have just been specified.
The discriminant retains quadratic vanishing on each branch. Its
zero scheme in the node has coordinate ring
\[
 R_{\rm node}/(u^2+v^2)
 =\mathbb C\{1,u,v,u^2\},\qquad
 v^2=-u^2,\quad u^3=v^3=0.
\]
These relations span the quotient, and independence follows by
homogeneous degree: degrees zero and one have no extra relation,
while in degree two precisely the one relation $u^2+v^2$ is imposed.
It consequently has length four. The conductor quotient
$R_{\rm node}/\mathfrak c=\mathbb C$ has length one. Their support
is the same origin because the radical of $(u^2+v^2)$ in the node
is $(u,v)$, as the displayed nilpotence proves. These are the exact
scheme and module relationships, including the multiplicity retained
by the discriminant.

\subsection{Exact induction to the connected general linear group}

We now turn the $K$-representations into explicitly defined rational
$G_{\rm node}$-representations. Put
\[
 L=\mathbb C[G_{\rm node}]
   =\mathbb C[p,q,s,t,(pt-qs)^{-1}].
\]
For any rational $K$-module $M$ define
\[
 \mathcal I(M)=(L\otimes_{\mathbb C}M)^K,
\]
where the diagonal action on a pure tensor is
$k\cdot(f\otimes m)=(g\mapsto f(gk))\otimes(k\cdot m)$.
The left $G_{\rm node}$-action is
$h\cdot(f\otimes m)=(g\mapsto f(h^{-1}g))\otimes m$.
The two actions commute by direct substitution; hence the invariant
subspace is a rational $G_{\rm node}$-module. When $M$ is a
$K$-algebra, $\mathcal I(M)$ is an algebra with componentwise
multiplication. These formulas define the invariant rings directly;
no assertion that their spectra are geometric quotients is needed.

We prove exactness rather than assume it. A rational torus module
is a sum of its integer weight spaces: expanding its coaction in the
independent Laurent monomials $\lambda^i\mu^j$ and comparing the
coaction identity proves that each coefficient is a vector of that
weight and that each vector has a finite such expansion. For a
surjective rational $T_{\rm node}$-module map and an invariant vector
in its target, choose a lift and take its weight-zero component;
the latter is an invariant lift. The normalizing swap preserves
weight zero. If the target vector is also swap invariant, average
that lift with its swap, dividing by two. This is a $K$-invariant
lift. Injections and kernels commute with invariants simply by
their definitions. Thus taking $K$-invariants is exact. Tensoring
an exact sequence of complex vector spaces with $L$ is exact, by
choosing bases or a vector-space splitting. This proves exactness
of $\mathcal I$ in full.

Apply it to the two proved sequences. With
$\mathcal A=\mathcal I(R_{\rm node})$ and
$\widetilde{\mathcal A}=\mathcal I(\widetilde R)$, we obtain
\begin{align}
 0&\longrightarrow\mathcal A\longrightarrow
 \widetilde{\mathcal A}\xrightarrow{\mathcal I(\epsilon)}
 \mathcal I(\mathbb C_{\rm sgn})\longrightarrow0,
 \label{eq:invariant-induced-ring}\\
 0&\longrightarrow\mathcal I(\mathbb C_{\det^{-1}})
 \longrightarrow\mathcal I(\Omega^1_{R_{\rm node}})
 \longrightarrow\mathcal I(\Omega^1_{\widetilde R})
 \longrightarrow0.
 \label{eq:invariant-induced-differential}
\end{align}
The maps are the original inclusion, difference of constants and
differential comparison applied coefficient by coefficient in $L$.
The first sequence is a sequence of $\mathcal A$-modules; its
quotient carries the action induced from evaluation at the node
origin. We have not identified the last two modules of
\eqref{eq:invariant-induced-differential} with absolute K\"ahler
differentials of the invariant rings. Such an identification is
unnecessary for the displayed exact transport.

The quotient and the kernel have explicit nonzero elements and
finite-dimensional $G_{\rm node}$-subrepresentations. Let
$J=\operatorname{diag}(1,-1)$ and define
\begin{equation}\label{eq:invariant-adjoint}
 Z(g)=gJg^{-1}
 =\frac1{pt-qs}
 \begin{pmatrix}pt+qs&-2pq\\2st&-(pt+qs)\end{pmatrix}.
\end{equation}
For a diagonal $k$ one has $Z(gk)=Z(g)$; for $k=\sigma$ one has
$Z(g\sigma)=-Z(g)$, since $\sigma J\sigma=-J$.
Consequently each entry of $Z(g)$ tensored with the fixed generator
of $\mathbb C_{\rm sgn}$ belongs to
$\mathcal I(\mathbb C_{\rm sgn})$. A lift through
\eqref{eq:invariant-induced-ring} is obtained by tensoring the same
entry with $(1/2,-1/2)\in\widetilde R$, whose difference of constants
is one and whose $K$-character is sign. The diagonal entry is nonzero
already at $g=I$.

The three coefficient functions $(pt+qs)/(pt-qs)$,
$-2pq/(pt-qs)$ and $2st/(pt-qs)$ are linearly independent. If a
linear combination were zero, multiplication by $pt-qs$ would give
a polynomial linear combination of $pt+qs,pq,st$; comparison of
their four distinct monomials forces its three coefficients zero.
They span a $G_{\rm node}$-stable three-dimensional module, because
$Z(h^{-1}g)=h^{-1}Z(g)h$. More exactly, with
$\mathfrak{sl}_2=\{X:\operatorname{tr}X=0\}$, the map
\[
 \mathfrak{sl}_2\longrightarrow\mathcal I(\mathbb C_{\rm sgn}),
 \qquad X\longmapsto\bigl(g\mapsto\operatorname{tr}(XZ(g))\bigr)
                 \otimes1
\]
is injective by the established coefficient independence and the
nondegenerate trace pairing on trace-zero two-by-two matrices.
It is equivariant for the adjoint action, since
\[
 \operatorname{tr}(X Z(h^{-1}g))
 =\operatorname{tr}(hXh^{-1}Z(g)).
\]
If desired its irreducibility can be checked without a classification
theorem: under diagonal conjugation the matrices $E_{12}$,
$\operatorname{diag}(1,-1)$ and $E_{21}$ have three distinct weights.
A nonzero invariant subspace therefore contains one of these weight
vectors, by coefficient projection. Differentiating conjugation by
$I+tE_{12}$ and $I+tE_{21}$ gives brackets; the identities
$[E_{12},E_{21}]=J$, $[J,E_{12}]=2E_{12}$ and
$[J,E_{21}]=-2E_{21}$ then produce all three basis vectors.

For the torsion kernel use the regular function $\det g=pt-qs$.
The tensor
\begin{equation}\label{eq:invariant-induced-torsion}
 e(g)=(\det g)\,\theta
 \in\mathcal I(\mathbb C_{\det^{-1}})
\end{equation}
is invariant because right translation multiplies its first factor
by $\det k$ and its second factor by $(\det k)^{-1}$.
It is nonzero, and the left action sends it to
$(\det h)^{-1}e$. Thus the exact connected-group transport retains
a determinant-inverse line in degree two as well as the adjoint
subrepresentation in the degree-zero conductor quotient. Degrees
refer to the node variables; $L$ has degree zero. The original
annihilator information is retained by the module maps: an element
of $\mathcal A$ whose coefficient functions vanish at the node origin
annihilates both displayed defect modules, because each original
defect is supported at that origin.

\subsection{The exact coordinate-ring obstruction interface}

For any two closed $G$-stable subsets $X,Y$ of a finite-dimensional
representation $W$, their polynomial ideals and coordinate rings
are defined by
$I(X)=\{f\in\mathbb C[W]:f|_X=0\}$ and
$\mathbb C[X]=\mathbb C[W]/I(X)$. Inclusion $X\subseteq Y$ gives
$I(Y)\subseteq I(X)$, and the entire equivariant restriction map is
\[
 \mathbb C[Y]\twoheadrightarrow\mathbb C[X],\qquad
 [f]_{I(Y)}\longmapsto[f]_{I(X)},\qquad
 \ker=I(X)/I(Y).
\]
Conversely, if a polynomial $f\in I(Y)$ has $f(x)\ne0$ at some
$x\in X$, that evaluation proves $x\notin Y$ and excludes the
inclusion. A finite-dimensional $G$-submodule containing that
polynomial may be used to record its representation content; it
is finite-dimensional because substitution in a polynomial of
bounded degree preserves that degree bound. This is the precise
polynomial-map interface used by an orbit-closure obstruction.

For the squarefree source cubic, the full ideal is zero by
\eqref{eq:invariant-open-orbit}, and its degreewise characters are
exactly \eqref{eq:invariant-cubic-character}. Thus this particular
orbit closure supplies no nonzero polynomial vanishing on it.
The saturation calculation proves exactly the same zero-kernel
statement for the whole original coefficient slice. The conductor
transport instead supplies the nonzero kernels and cokernels in
\eqref{eq:invariant-induced-ring}--\eqref{eq:invariant-induced-differential},
with explicit module elements and representations. Those sequences
are actual correspondences with connected-group representation
data; they do not assert that these modules are the coordinate ring
of a permanent or determinant orbit closure. The proved common
edge between the original node and the binary-cubic representation
is \eqref{eq:invariant-node-to-cubic}, with its exact kernel,
discriminant pullback, degrees, and torus homomorphisms. Every claim
in this fixed-dimensional transport is thus attached to a specified
map. No assertion about the complexity of an unbounded family or
about a nonstandard quantum-group positivity conjecture has been
substituted for one of these calculations.
